\documentclass[10.5pt,compsoc]{CjC}
\usepackage{CJKutf8}
\usepackage{amsmath}
\usepackage{booktabs}
\usepackage{graphicx}
\usepackage{url}
\usepackage[noadjust]{cite}
\usepackage{captionhack}
\setlength{\parindent}{2em}
\setcounter{page}{1}

\begin{document}
\begin{CJK*}{UTF8}{gbsn}

\twocolumn[
\begin{center}
{\zihao{2}\bfseries 基于Apache Arrow与CUDA内存模式的TPC-H Q5查询实验\par}
\vspace{4mm}
{\zihao{4}徐子桓\quad 2025104082\par}
\vspace{4mm}
\begin{minipage}{0.94\textwidth}
\zihao{5}
\sloppy
\textbf{摘\quad 要}\quad
本文围绕内存数据库中的CPU--GPU协同查询处理，实现了一个固定的TPC-H
Q5物理执行器。实验保留Q5需要的六张表和字段，使用Apache Arrow IPC作为
可复现的数据交换格式，并比较Arrow CPU算子、专用CPU实现、三种CUDA内存
模式以及cuDF算子库。三种CUDA模式分别是显式PCIe复制、managed memory
预取和mapped pinned host memory直接访问。项目先用tiny数据检查查询逻辑，
再在官方TPC-H SF1数据上验证结果并测量时间。当前实现没有做完整SQL解析和
代价优化，GPU也只负责最后的lineitem扫描，因此结论只适用于本项目的固定
查询和SF1规模。实验中还发现，原实现逐行截断收入会偏离官方答案，因此本文
改为保留四位小数累加，最后显示时再舍入。本轮六种实现得到相同结果；在三次
重复的内部计时中，8线程专用CPU为22.86 ms，三种CUDA路径为238.70至
323.24 ms。冷进程墙钟时间则受到输入格式影响，Arrow IPC路径明显少于C++
文本装载路径。

\textbf{关键词}\quad 内存数据库；TPC-H Q5；Apache Arrow；CUDA；统一虚拟寻址；cuDF
\end{minipage}
\vspace{5mm}

\begin{minipage}{0.94\textwidth}
\zihao{5}
\textbf{Abstract}\quad
This report studies CPU--GPU query processing with a fixed implementation of
TPC-H Query 5. Apache Arrow IPC is used as the reproducible interchange format.
The experiment compares an Arrow CPU baseline, a specialized CPU executor,
three CUDA memory modes, and a cuDF operator baseline. The current prototype
does not implement a general SQL optimizer, and the GPU only scans and
aggregates the lineitem table. Therefore, the conclusions are limited to the
implemented query and the SF1 data set.

\textbf{Keywords}\quad in-memory database; TPC-H Q5; Apache Arrow; CUDA; UVA; cuDF
\end{minipage}
\vspace{5mm}
\end{center}
]

\zihao{5}

\section{引言}

GPU具有较高的并行计算能力和显存带宽，但数据库查询不只包含计算，还包含
数据装载、连接准备、主机与设备之间的数据移动等步骤。只比较一个CUDA核函数
容易得到不完整的结论。因此，本实验选用TPC-H Q5，把CPU列式算子、专用CPU
执行器、不同GPU内存模式和GPU算子库放在同一查询语义下进行对比。

本项目不是通用数据库。它没有SQL解析器、逻辑优化器和事务模块，而是把Q5
的连接顺序和过滤过程直接写成物理执行计划。这样实现范围较小，适合课程实验，
但也意味着结果不能直接推广到其他查询。

\section{查询与数据表示}

Q5包含六张表。region和nation用于确定地区与国家，customer和orders用于
确定客户与订单，supplier和lineitem用于供应商匹配和收入计算。查询选择
指定地区及一年日期范围内的订单，要求客户与供应商属于同一国家，最后按国家聚合
\begin{equation}
  \mathrm{revenue}=\sum l\_extendedprice\times(1-l\_discount).
\end{equation}

项目只保留查询需要的列。字符串使用字典编码，日期转换为从Unix epoch开始
的天数。价格以分为单位保存，折扣以百分之一为单位保存。单行收入用整数表示
为
\begin{equation}
  r_{10^4}=p_{cent}\times(100-d_{0.01}),
\end{equation}
先按$10^4$尺度求和，输出时再舍入到两位小数。这一处理避免原版本逐行截断
不足一分部分的问题。

数据准备脚本把TPC-H的六个\texttt{.tbl}文件转换为Arrow IPC文件，同时写入
schema、行数和SHA256校验信息。Arrow数据集作为PyArrow与cuDF基线的共同
输入。专用C++/CUDA路径目前仍从\texttt{.tbl}构造连续数组，这是本实验尚未
统一的地方。

\section{查询执行方法}

\subsection{CPU过滤传播}

执行器先在CPU端找到目标region中的nation，然后构造supplier、customer和
order到nation的三个直接索引数组，以下分别记为$S[s]$、$C[c]$和$O[o]$。
日期或地区条件不满足的
位置写为$-1$。最后扫描lineitem，只要订单国家与供应商国家相同，就把收入
加入对应国家。该方法避免物化六表连接，但依赖TPC-H键值比较稠密的特点。

专用CPU版本把lineitem划分为连续区间，每个线程使用本地国家聚合数组，结束
后再归并。Arrow CPU版本使用PyArrow的filter、join、group by和sort算子，
作为通用列式算子库对照。

\subsection{CUDA内存模式}

三个手写CUDA后端复用CPU生成的过滤传播数组，并执行同一个lineitem聚合核：

\begin{itemize}
  \item \texttt{gpu-copy}使用device buffer和显式复制；
  \item \texttt{gpu-managed}使用统一内存并在执行前预取；
  \item \texttt{gpu-mapped}使用映射锁页内存，GPU经PCIe读取主机数据。
\end{itemize}

这三个模式比较的是输入数组到GPU的存放和访问方式，并不是三种不同的连接
算法。cuDF版本则通过通用GPU DataFrame的merge、filter、groupby和sort算子
执行Q5，用来观察专用核函数与算子库之间的差别。

\section{实验设计}

实验分为两级。tiny数据只有6条lineitem，用来检查日期、地区、同国家条件和
收入公式。正式数据使用TPC-H V3.0.1 dbgen生成的SF1。每个可用后端先预热，
再独立进程重复3次并保存原始CSV、环境信息和结果哈希。服务器包含2路AMD
EPYC 9654和NVIDIA GeForce RTX 4090，C++路径使用8个线程，软件环境为
CUDA 12.6、PyArrow 23.0.1和cuDF 26.06.00。正确性要求不同后端输出的国家
顺序和$10^4$尺度收入相同，并与官方\texttt{q5.out}的两位小数结果一致。

为了缩短课程项目周期，正式实验只做SF1，没有继续做SF10；也没有做查询并发、
NUMA绑定和完整Nsight分析。因此，本文主要回答“这些实现能否正确运行、主要
时间花在哪里”，不尝试给出普遍的性能模型。

\section{实验结果}

\subsection{正确性}

tiny数据的预期结果为JAPAN 190.00、INDIA 90.00。修正decimal语义后，CPU
与三种CUDA模式在$10^4$尺度上的结果哈希一致。

SF1上六种实现的结果哈希均为\texttt{542abf4003633c7c}。表\ref{tab:oracle}
列出的五行与官方\texttt{q5.out}逐行一致，核对过程没有把字符串转换为浮点数。

\begin{table}[htbp]
\centering
\caption{SF1 Q5结果与官方答案核对}
\label{tab:oracle}
\small
\begin{tabular}{lr}
\toprule
国家 & 收入\\
\midrule
INDONESIA & 55502041.17\\
VIETNAM & 55295087.00\\
CHINA & 53724494.26\\
INDIA & 52035512.00\\
JAPAN & 45410175.70\\
\bottomrule
\end{tabular}
\end{table}

\subsection{性能观察}

表\ref{tab:perf}给出一次预热、3次正式运行的中位数。内部\texttt{total}
由各后端报告，冷进程墙钟由外部控制器测量。专用CPU在数据已装入内存后的
查询主体最短。三种CUDA路径都没有在SF1一次性查询中抵消初始化、分配和内存
管理成本。

\begin{table}[htbp]
\centering
\caption{SF1时间中位数（ms）}
\label{tab:perf}
\small
\begin{tabular}{lrr}
\toprule
后端 & 内部total & 冷进程墙钟\\
\midrule
CPU (8线程) & 22.859 & 5373.400\\
gpu-copy & 238.698 & 5597.830\\
gpu-managed & 270.005 & 5654.902\\
gpu-mapped & 323.239 & 5717.684\\
PyArrow & 1069.345 & 1279.315\\
cuDF & 1907.409 & 2347.973\\
\bottomrule
\end{tabular}
\end{table}

更细的CUDA计时能说明内存模式差异。gpu-copy的H2D中位数为6.577 ms，
kernel为0.230 ms；managed的预取阶段为6.299 ms，kernel为0.216 ms；
mapped的显式H2D只有0.016 ms，但kernel增加到2.703 ms。这说明mapped并不是
取消PCIe访问，而是把传输从显式复制变成核函数运行期间的远程读取。

两种时间列不能混为同一种排名。C++/CUDA的内部计时在\texttt{.tbl}已读入
内存后开始，而PyArrow/cuDF把Arrow IPC读取和对象构造包含在内部计时中；
外部墙钟统一包含子进程启动和装载。Arrow墙钟较短主要反映IPC列式输入比当前
C++文本解析更合适，不能单独证明通用Arrow算子的查询执行快于专用CPU。

\section{讨论与不足}

本实验有几处比较明显的不足。第一，专用C++路径和Arrow/cuDF路径还没有完全
统一数据装载口径；前者读取文本并构造数组，后者读取Arrow IPC。第二，GPU只
做lineitem扫描，不是完整的GPU查询计划。第三，三种CUDA模式每次运行都重新
分配内存，无法代表常驻数据库缓冲池。第四，只测试SF1，数据量不足以说明更大
规模下的交叉点。第五，聚合使用全局原子加，没有实现block内局部归约。

这些问题不影响本次“实现并比较不同内存模式”的基本目标，但会限制性能结论。
后续最直接的改进是让所有后端从同一Arrow批次装载、复用GPU缓冲区，再增加
SF10和多查询实验。

\section{结论}

本文完成了一个可以复现的TPC-H Q5 CPU--GPU实验原型。实现表明，GPU内存
模式的差别不仅是API差别，还会改变PCIe传输发生的位置和核函数的数据访问
代价。对于当前一次性执行的SF1查询，专用CPU的查询主体最短，GPU固定开销
没有被六百万行lineitem摊薄；但统一墙钟又显示文本装载会反过来成为主要成本。
项目同时通过官方答案发现并修复了原有decimal截断问题。由于系统
仍是固定查询执行器、实验规模也较小，本文把更大规模和完整GPU查询计划留作
后续工作。

\begin{thebibliography}{9}
\zihao{5-}
\bibitem{tpch}
Transaction Processing Performance Council. TPC Benchmark H Standard
Specification Revision 3.0.1[EB/OL]. \url{https://www.tpc.org/tpch/}.

\bibitem{arrow}
Apache Arrow Project. Apache Arrow Documentation[EB/OL].
\url{https://arrow.apache.org/docs/}.

\bibitem{cuda}
NVIDIA. CUDA C++ Programming Guide[EB/OL].
\url{https://docs.nvidia.com/cuda/cuda-c-programming-guide/}.

\bibitem{cudf}
RAPIDS. cuDF Documentation[EB/OL].
\url{https://docs.rapids.ai/api/cudf/stable/}.
\end{thebibliography}

\end{CJK*}
\end{document}
